%=====================================
%   20250723.tex
%   20250723 
%   toshi2019
%=====================================
% ------------------------
% documentclass
% ------------------------
\documentclass[leqno, 10pt, a4paper, dvipdfmx]{jsarticle}
\title{ゼミノート（07月23日分）}
\author{Toshi2019}
\date{\today}

% ------------------------
% usepackage
% ------------------------
\usepackage{algorithm}
\usepackage{algorithmic}
\usepackage{amscd}
\usepackage{amsfonts}
\usepackage{amsmath}
\usepackage[psamsfonts]{amssymb}
\usepackage{amsthm}
\usepackage{ascmac}
\usepackage{bm}
\usepackage{color}
\usepackage{enumerate}
\usepackage{fancybox}
\usepackage[stable]{footmisc}
\usepackage{graphicx}
\usepackage{listings}
\usepackage{mathrsfs}
\usepackage{mathtools}
\usepackage{otf}
\usepackage{pifont}
\usepackage{proof}
\usepackage{subfigure}
\usepackage{tikz}
\usepackage{verbatim}
\usepackage[all]{xy}

\usetikzlibrary{cd}
\usetikzlibrary{arrows.meta}
\usetikzlibrary{calc}

% ================================
% パッケージを追加する場合のスペース 
\usepackage[dvipdfmx]{hyperref}
\usepackage{xcolor}
\definecolor{darkgreen}{rgb}{0,0.45,0} 
\definecolor{darkred}{rgb}{0.75,0,0}
\definecolor{darkblue}{rgb}{0,0,0.6} 
\hypersetup{
    colorlinks=true,
    citecolor=darkgreen,
    linkcolor=darkred,
    urlcolor=darkblue,
}
\usepackage{pxjahyper}

%\usepackage[mathscr]{euscript} % mathscr のフォントを変更

%\usepackage{mmasym}

%=================================


% --------------------------
% theoremstyle
% --------------------------
\theoremstyle{definition}

% --------------------------
% newtheoem
% --------------------------

% 日本語で定理, 命題, 証明などを番号付きで用いるためのコマンドです. 
% If you want to use theorem environment in Japanece, 
% you can use these code. 
% Attention!
% All theorem enivironment numbers depend on 
% only section numbers.
\newtheorem{Axiom}{公理}[section]
\newtheorem{Definition}[Axiom]{定義}
\newtheorem{Theorem}[Axiom]{定理}
\newtheorem{Proposition}[Axiom]{命題}
\newtheorem{Lemma}[Axiom]{補題}
\newtheorem{Corollary}[Axiom]{系}
\newtheorem{Example}[Axiom]{例}
\newtheorem{Claim}[Axiom]{主張}
\newtheorem{Property}[Axiom]{性質}
\newtheorem{Attention}[Axiom]{注意}
\newtheorem{Question}[Axiom]{問}
\newtheorem{Problem}[Axiom]{問題}
\newtheorem{Consideration}[Axiom]{考察}
\newtheorem{Alert}[Axiom]{警告}
\newtheorem{Fact}[Axiom]{事実}


% 日本語で定理, 命題, 証明などを番号なしで用いるためのコマンドです. 
% If you want to use theorem environment with no number in Japanese, You can use these code.
\newtheorem*{Axiom*}{公理}
\newtheorem*{Definition*}{定義}
\newtheorem*{Theorem*}{定理}
\newtheorem*{Proposition*}{命題}
\newtheorem*{Lemma*}{補題}
\newtheorem*{Example*}{例}
\newtheorem*{Corollary*}{系}
\newtheorem*{Claim*}{主張}
\newtheorem*{Property*}{性質}
\newtheorem*{Attention*}{注意}
\newtheorem*{Question*}{問}
\newtheorem*{Problem*}{問題}
\newtheorem*{Consideration*}{考察}
\newtheorem*{Alert*}{警告}
\newtheorem{Fact*}{事実}


% 英語で定理, 命題, 証明などを番号付きで用いるためのコマンドです. 
%　If you want to use theorem environment in English, You can use these code.
%all theorem enivironment number depend on only section number.
\newtheorem{Axiom+}{Axiom}[section]
\newtheorem{Definition+}[Axiom+]{Definition}
\newtheorem{Theorem+}[Axiom+]{Theorem}
\newtheorem{Proposition+}[Axiom+]{Proposition}
\newtheorem{Lemma+}[Axiom+]{Lemma}
\newtheorem{Example+}[Axiom+]{Example}
\newtheorem{Corollary+}[Axiom+]{Corollary}
\newtheorem{Claim+}[Axiom+]{Claim}
\newtheorem{Property+}[Axiom+]{Property}
\newtheorem{Attention+}[Axiom+]{Attention}
\newtheorem{Question+}[Axiom+]{Question}
\newtheorem{Problem+}[Axiom+]{Problem}
\newtheorem{Consideration+}[Axiom+]{Consideration}
\newtheorem{Alert+}{Alert}
\newtheorem{Fact+}[Axiom+]{Fact}
\newtheorem{Remark+}[Axiom+]{Remark}

% ----------------------------
% commmand
% ----------------------------
% 執筆に便利なコマンド集です. 
% コマンドを追加する場合は下のスペースへ. 

% 集合の記号 (黒板文字)
\newcommand{\NN}{\mathbb{N}}
\newcommand{\ZZ}{\mathbb{Z}}
\newcommand{\QQ}{\mathbb{Q}}
\newcommand{\RR}{\mathbb{R}}
\newcommand{\CC}{\mathbb{C}}
\newcommand{\PP}{\mathbb{P}}
\newcommand{\KK}{\mathbb{K}}


% 集合の記号 (太文字)
\newcommand{\nn}{\mathbf{N}}
\newcommand{\zz}{\mathbf{Z}}
\newcommand{\qq}{\mathbf{Q}}
\newcommand{\rr}{\mathbf{R}}
\newcommand{\cc}{\mathbf{C}}
\newcommand{\pp}{\mathbf{P}}
\newcommand{\kk}{\mathbf{k}}

% 特殊な写像の記号
\newcommand{\ev}{\mathop{\mathrm{ev}}\nolimits} % 値写像
\newcommand{\pr}{\mathop{\mathrm{pr}}\nolimits} % 射影

% スクリプト体にするコマンド
%   例えば　{\mcal C} のように用いる
\newcommand{\mcal}{\mathcal}

% 花文字にするコマンド 
%   例えば　{\h C} のように用いる
\newcommand{\h}{\mathscr}

% ヒルベルト空間などの記号
\newcommand{\F}{\mcal{F}}
\newcommand{\X}{\mcal{X}}
\newcommand{\Y}{\mcal{Y}}
\newcommand{\Hil}{\mcal{H}}
\newcommand{\RKHS}{\Hil_{k}}
\newcommand{\Loss}{\mcal{L}_{D}}
\newcommand{\MLsp}{(\X, \Y, D, \Hil, \Loss)}

% 偏微分作用素の記号
\newcommand{\p}{\partial}

% 角カッコの記号 (内積は下にマクロがあります)
\newcommand{\lan}{\langle}
\newcommand{\ran}{\rangle}



% 圏の記号など
\newcommand{\Set}{{\bf Set}}
\newcommand{\Vect}{{\bf Vect}}
\newcommand{\FDVect}{{\bf FDVect}}
%\newcommand{\Ring}{{\bf Ring}}
\newcommand{\Ab}{{\bf Ab}}
\newcommand{\Mod}{\mathop{\mathrm{Mod}}\nolimits}
\newcommand{\Modf}{\mathop{\mathrm{Mod}^\mathrm{f}}\nolimits}
\newcommand{\CGA}{{\bf CGA}}
\newcommand{\GVect}{{\bf GVect}}
\newcommand{\Lie}{{\bf Lie}}
\newcommand{\dLie}{{\bf Liec}}



% 射の集合など
\newcommand{\Map}{\mathop{\mathrm{Map}}\nolimits} % 写像の集合
\newcommand{\Hom}{\mathop{\mathrm{Hom}}\nolimits} % 射集合
\newcommand{\End}{\mathop{\mathrm{End}}\nolimits} % 自己準同型の集合
\newcommand{\Aut}{\mathop{\mathrm{Aut}}\nolimits} % 自己同型の集合
\newcommand{\Mor}{\mathop{\mathrm{Mor}}\nolimits} % 射集合
\newcommand{\Ker}{\mathop{\mathrm{Ker}}\nolimits} % 核
\newcommand{\Img}{\mathop{\mathrm{Im}}\nolimits} % 像
\newcommand{\Cok}{\mathop{\mathrm{Coker}}\nolimits} % 余核
\newcommand{\Cim}{\mathop{\mathrm{Coim}}\nolimits} % 余像

% その他便利なコマンド
\newcommand{\dip}{\displaystyle} % 本文中で数式モード
\newcommand{\e}{\varepsilon} % イプシロン
\newcommand{\dl}{\delta} % デルタ
\newcommand{\pphi}{\varphi} % ファイ
\newcommand{\ti}{\tilde} % チルダ
\newcommand{\pal}{\parallel} % 平行
\newcommand{\op}{{\rm op}} % 双対を取る記号
\newcommand{\lcm}{\mathop{\mathrm{lcm}}\nolimits} % 最小公倍数の記号
\newcommand{\Probsp}{(\Omega, \F, \P)} 
\newcommand{\argmax}{\mathop{\rm arg~max}\limits}
\newcommand{\argmin}{\mathop{\rm arg~min}\limits}





% ================================
% コマンドを追加する場合のスペース 
%\newcommand{\OO}{\mcal{O}}



\makeatletter
\renewenvironment{proof}[1][\proofname]{\par
  \pushQED{\qed}%
  \normalfont \topsep6\p@\@plus6\p@\relax
  \trivlist
  \item[\hskip\labelsep
%        \itshape
         \bfseries
%    #1\@addpunct{.}]\ignorespaces
    {#1}]\ignorespaces
}{%
  \popQED\endtrivlist\@endpefalse
}
\makeatother

\renewcommand{\proofname}{証明.}



%\renewcommand\proofname{\bf 証明} % 証明
\numberwithin{equation}{section}
\newcommand{\cTop}{\textsf{Top}}
%\newcommand{\cOpen}{\textsf{Open}}
\newcommand{\Op}{\mathop{\textsf{Op}}\nolimits}
\newcommand{\Ob}{\mathop{\textrm{Ob}}\nolimits}
\newcommand{\id}{\mathop{\mathrm{id}}\nolimits}
\newcommand{\pt}{\mathop{\mathrm{pt}}\nolimits}
\newcommand{\res}{\mathop{\rho}\nolimits}
\newcommand{\A}{\mcal{A}}
\newcommand{\B}{\mcal{B}}
\newcommand{\C}{\mcal{C}}
\newcommand{\D}{\mcal{D}}
\newcommand{\E}{\mcal{E}}
\newcommand{\G}{\mcal{G}}
%\newcommand{\H}{\mcal{H}}
\newcommand{\I}{\mcal{I}}
\newcommand{\J}{\mcal{J}}
\newcommand{\OO}{\mcal{O}}
\newcommand{\Ring}{\mathop{\textsf{Ring}}\nolimits}
\newcommand{\cAb}{\mathop{\textsf{Ab}}\nolimits}
%\newcommand{\Ker}{\mathop{\mathrm{Ker}}\nolimits}
\newcommand{\im}{\mathop{\mathrm{Im}}\nolimits}
\newcommand{\Coker}{\mathop{\mathrm{Coker}}\nolimits}
\newcommand{\Coim}{\mathop{\mathrm{Coim}}\nolimits}
\newcommand{\rank}{\mathop{\mathrm{rank}}\nolimits}
\newcommand{\Ht}{\mathop{\mathrm{Ht}}\nolimits}
\newcommand{\supp}{\mathop{\mathrm{supp}}\nolimits}
\newcommand{\colim}{\mathop{\mathrm{colim}}}
\newcommand{\Tor}{\mathop{\mathrm{Tor}}\nolimits}

\newcommand{\cat}{\mathscr{C}}

%筆記体
\newcommand{\cA}{\mcal{A}}
\newcommand{\cB}{\mcal{B}}
\newcommand{\cC}{\mcal{C}}
\newcommand{\cD}{\mcal{D}}
\newcommand{\cE}{\mcal{E}}
\newcommand{\cF}{\mcal{F}}
\newcommand{\cG}{\mcal{G}}
\newcommand{\cH}{\mcal{H}}
\newcommand{\cI}{\mcal{I}}
\newcommand{\cJ}{\mcal{J}}
\newcommand{\cK}{\mcal{K}}
\newcommand{\cL}{\mcal{L}}
\newcommand{\cM}{\mcal{M}}
\newcommand{\cN}{\mcal{N}}
\newcommand{\cO}{\mcal{O}}
\newcommand{\cP}{\mcal{P}}
\newcommand{\cQ}{\mcal{Q}}
\newcommand{\cR}{\mcal{R}}
\newcommand{\cS}{\mcal{S}}
\newcommand{\cT}{\mcal{T}}
\newcommand{\cU}{\mcal{U}}
\newcommand{\cV}{\mcal{V}}
\newcommand{\cW}{\mcal{W}}
\newcommand{\cX}{\mcal{X}}
\newcommand{\cY}{\mcal{Y}}
\newcommand{\cZ}{\mcal{Z}}


\newcommand{\scA}{\mathscr{A}}
\newcommand{\scB}{\mathscr{B}}
\newcommand{\scC}{\mathscr{C}}
\newcommand{\scD}{\mathscr{D}}
\newcommand{\scE}{\mathscr{E}}
\newcommand{\scF}{\mathscr{F}}
\newcommand{\scN}{\mathscr{N}}
\newcommand{\scO}{\mathscr{O}}
\newcommand{\scV}{\mathscr{V}}
\newcommand{\scU}{\mathscr{U}}


\newcommand{\ibA}{\mathop{\text{\textit{\textbf{A}}}}}
\newcommand{\ibB}{\mathop{\text{\textit{\textbf{B}}}}}
\newcommand{\ibC}{\mathop{\text{\textit{\textbf{C}}}}}
\newcommand{\ibD}{\mathop{\text{\textit{\textbf{D}}}}}
\newcommand{\ibE}{\mathop{\text{\textit{\textbf{E}}}}}
\newcommand{\ibF}{\mathop{\text{\textit{\textbf{F}}}}}
\newcommand{\ibG}{\mathop{\text{\textit{\textbf{G}}}}}
\newcommand{\ibH}{\mathop{\text{\textit{\textbf{H}}}}}
\newcommand{\ibI}{\mathop{\text{\textit{\textbf{I}}}}}
\newcommand{\ibJ}{\mathop{\text{\textit{\textbf{J}}}}}
\newcommand{\ibK}{\mathop{\text{\textit{\textbf{K}}}}}
\newcommand{\ibL}{\mathop{\text{\textit{\textbf{L}}}}}
\newcommand{\ibM}{\mathop{\text{\textit{\textbf{M}}}}}
\newcommand{\ibN}{\mathop{\text{\textit{\textbf{N}}}}}
\newcommand{\ibO}{\mathop{\text{\textit{\textbf{O}}}}}
\newcommand{\ibP}{\mathop{\text{\textit{\textbf{P}}}}}
\newcommand{\ibQ}{\mathop{\text{\textit{\textbf{Q}}}}}
\newcommand{\ibR}{\mathop{\text{\textit{\textbf{R}}}}}
\newcommand{\ibS}{\mathop{\text{\textit{\textbf{S}}}}}
\newcommand{\ibT}{\mathop{\text{\textit{\textbf{T}}}}}
\newcommand{\ibU}{\mathop{\text{\textit{\textbf{U}}}}}
\newcommand{\ibV}{\mathop{\text{\textit{\textbf{V}}}}}
\newcommand{\ibW}{\mathop{\text{\textit{\textbf{W}}}}}
\newcommand{\ibX}{\mathop{\text{\textit{\textbf{X}}}}}
\newcommand{\ibY}{\mathop{\text{\textit{\textbf{Y}}}}}
\newcommand{\ibZ}{\mathop{\text{\textit{\textbf{Z}}}}}

\newcommand{\ibx}{\mathop{\text{\textit{\textbf{x}}}}}

%\newcommand{\Comp}{\mathop{\mathrm{C}}\nolimits}
%\newcommand{\Komp}{\mathop{\mathrm{K}}\nolimits}
%\newcommand{\Domp}{\mathop{\mathsf{D}}\nolimits}%複体のホモトピー圏
%\newcommand{\Comp}{\mathrm{C}}
%\newcommand{\Komp}{\mathrm{K}}
%\newcommand{\Domp}{\mathsf{D}}%複体のホモトピー圏

\newcommand{\Comp}{\mathop{\mathrm{C}}\nolimits}
\newcommand{\Komp}{\mathop{\mathsf{K}}\nolimits}
\newcommand{\Domp}{\mathord{\mathsf{D}}}%\nolimits}
\newcommand{\Kompl}{\mathop{\mathsf{K}^\mathrm{+}}\nolimits}
\newcommand{\Kompu}{\mathop{\mathsf{K}^\mathrm{-}}\nolimits}
\newcommand{\Kompb}{\mathop{\mathsf{K}^\mathrm{b}}\nolimits}
\newcommand{\Dompl}{\mathop{\mathsf{D}^\mathrm{+}}\nolimits}
\newcommand{\Dompu}{\mathop{\mathsf{D}^\mathrm{-}}\nolimits}
\newcommand{\Dompb}{\mathop{\mathsf{D}^\mathrm{b}}\nolimits}
\newcommand{\Dbconic}{\mathop{\mathsf{D}^\mathrm{b}_{\rrp}}\nolimits}
\newcommand{\Dompbf}{\mathop{\mathsf{D}_\mathrm{f}^\mathrm{b}}\nolimits}
\newcommand{\Domplb}{\mathop{\mathsf{D}^\mathrm{lb}}\nolimits}




\newcommand{\CCat}{\Comp(\cat)}
\newcommand{\KCat}{\Komp(\cat)}
\newcommand{\DCat}{\Domp(\cat)}%圏Cの複体のホモトピー圏
%\newcommand{\HOM}{\mathop{\mathscr{H}\hspace{-2pt}om}\nolimits}%内部Hom
\newcommand{\HOM}{\mathop{\mathcal{H}\hspace{-0pt}om}\nolimits}%内部Hom
\newcommand{\RHOM}{\mathop{\mathrm{R}\hspace{-1.5pt}\HOM}\nolimits}

\newcommand{\muS}{\mathop{\mathrm{SS}}\nolimits}
\newcommand{\RG}{\mathop{\mathrm{R}\hspace{-0pt}\Gamma}\nolimits}
\newcommand{\RHom}{\mathop{\mathrm{R}\hspace{-1.5pt}\Hom}\nolimits}
\newcommand{\Rder}{\mathrm{R}}

\newcommand{\simar}{\mathrel{\overset{\sim}{\rightarrow}}}%同型右矢印
\newcommand{\simarr}{\mathrel{\overset{\sim}{\longrightarrow}}}%同型右矢印
\newcommand{\simra}{\mathrel{\overset{\sim}{\leftarrow}}}%同型左矢印
\newcommand{\simrra}{\mathrel{\overset{\sim}{\longleftarrow}}}%同型左矢印

\newcommand{\hocolim}{{\mathrm{hocolim}}}
\newcommand{\indlim}[1][]{\mathop{\varinjlim}\limits_{#1}}
\newcommand{\sindlim}[1][]{\smash{\mathop{\varinjlim}\limits_{#1}}\,}
\newcommand{\Pro}{\mathrm{Pro}}
\newcommand{\Ind}{\mathrm{Ind}}
\newcommand{\prolim}[1][]{\mathop{\varprojlim}\limits_{#1}}
\newcommand{\sprolim}[1][]{\smash{\mathop{\varprojlim}\limits_{#1}}\,}

\newcommand{\Sh}{\mathrm{Sh}}
\newcommand{\PSh}{\mathrm{PSh}}

\newcommand{\rmD}{\mathsf{D}}
\newcommand{\vD}{\mathbf{D}}

\newcommand{\Lloc}[1][]{\mathord{\mathcal{L}^1_{\mathrm{loc},{#1}}}}
\newcommand{\ori}{\mathord{\mathrm{or}}}
\newcommand{\Db}{\mathord{\mathscr{D}b}}

\newcommand{\codim}{\mathop{\mathrm{codim}}\nolimits}



\newcommand{\gld}{\mathop{\mathrm{gld}}\nolimits}
\newcommand{\wgld}{\mathop{\mathrm{wgld}}\nolimits}


\newcommand{\tens}[1][]{\mathbin{\otimes_{\raise1.5ex\hbox to-.1em{}{#1}}}}
\newcommand{\etens}{\mathbin{\boxtimes}}
\newcommand{\ltens}[1][]{\mathbin{\overset{\mathrm{L}}\tens}_{#1}}
\newcommand{\mtens}[1][]{\mathbin{\overset{\mathrm{\mu}}\tens}_{#1}}
\newcommand{\lltens}[1][]{\mathbin{{\mathop{\tens}\limits^{\mathrm{L}}_{#1}}}}
\newcommand{\lletens}[1][]{\mathbin{{\mathop{\boxtimes}\limits^{\mathrm{L}}_{#1}}}}
\newcommand{\letens}{\overset{\mathrm{L}}{\etens}}
\newcommand{\detens}{\underline{\etens}}
\newcommand{\ldetens}{\overset{\mathrm{L}}{\underline{\etens}}}
\newcommand{\dtens}[1][]{{\overset{\mathrm{L}}{\underline{\otimes}}}_{#1}}

\newcommand{\blk}{\mathord{\ \cdot\ }}
\newcommand{\mres}[2][]{{\left.{#1}\right\rvert}_{#2}}


%\newcommand{\hocolim}{{\mathrm{hocolim}}}
%\newcommand{\indlim}[1][]{\mathop{\varinjlim}\limits_{#1}}
%\newcommand{\sindlim}[1][]{\smash{\mathop{\varinjlim}\limits_{#1}}\,}
%\newcommand{\Pro}{\mathrm{Pro}}
%\newcommand{\Ind}{\mathrm{Ind}}
%\newcommand{\prolim}[1][]{\mathop{\varprojlim}\limits_{#1}}
%\newcommand{\sprolim}[1][]{\smash{\mathop{\varprojlim}\limits_{#1}}\,}
\newcommand{\proolim}[1][]{\mathop{\text{\rm``{$\varprojlim$}''}}\limits_{#1}}
\newcommand{\sproolim}[1][]{\smash{\mathop{\rm``{\varprojlim}''}\limits_{#1}}}
\newcommand{\inddlim}[1][]{\mathop{\text{\rm``{$\varinjlim$}''}}\limits_{#1}}
\newcommand{\sinddlim}[1][]{\smash{\mathop{\text{\rm``{$\varinjlim$}''}}\limits_{#1}}\,}
\newcommand{\ooplus}{\mathop{\text{\rm``{$\oplus$}''}}\limits}
\newcommand{\bbigsqcup}{\mathop{``\bigsqcup''}\limits}
\newcommand{\bsqcup}{\mathop{``\sqcup''}\limits}
\newcommand{\dsum}[1][]{\mathbin{\oplus_{#1}}}

\newcommand{\Fct}{\mathop{\mathsf{Fct}}\nolimits}

\newcommand{\pb}[1][]{\mathbin{\mathop{\times}\limits_{#1}}}
\newcommand{\dpb}[1][]{\mathbin{\mathop{\times}\limits_{#1}}}
\newcommand{\tpb}[1][]{\mathbin{\mathop{\times}\nolimits_{#1}}}





%================================================
% 自前の定理環境
%   https://mathlandscape.com/latex-amsthm/
% を参考にした
\newtheoremstyle{mystyle}%   % スタイル名
    {5pt}%                   % 上部スペース
    {5pt}%                   % 下部スペース
    {}%              % 本文フォント
    {}%                  % 1行目のインデント量
    {\bfseries}%                      % 見出しフォント
    {.}%                     % 見出し後の句読点
    {12pt}%                     % 見出し後のスペース
    {\thmname{#1}\thmnumber{ #2}\thmnote{{\hspace{2pt}\normalfont (#3)}}}% % 見出しの書式

\theoremstyle{mystyle}
\newtheorem{AXM}{公理}[section]
\newtheorem{DFN}[AXM]{定義}
\newtheorem{THM}[AXM]{定理}
\newtheorem*{THM*}{定理}
\newtheorem{PRP}[AXM]{命題}
\newtheorem{LMM}[AXM]{補題}
\newtheorem{CRL}[AXM]{系}
\newtheorem{EG}[AXM]{例}
\newtheorem*{EG*}{例}
\newtheorem{RMK}[AXM]{注意}
\newtheorem{CNV}[AXM]{約束}
\newtheorem{CMT}[AXM]{コメント}
\newtheorem*{CMT*}{コメント}
\newtheorem{NTN}[AXM]{記号}

% 定理環境ここまで
%====================================================

\usepackage{framed}
\definecolor{lightgray}{rgb}{0.75,0.75,0.75}
\renewenvironment{leftbar}{%
  \def\FrameCommand{\textcolor{lightgray}{\vrule width 4pt} \hspace{10pt}}% 
  \MakeFramed {\advance\hsize-\width \FrameRestore}}%
{\endMakeFramed}
\newenvironment{redleftbar}{%
  \def\FrameCommand{\textcolor{lightgray}{\vrule width 1pt} \hspace{10pt}}% 
  \MakeFramed {\advance\hsize-\width \FrameRestore}}%
 {\endMakeFramed}



\renewcommand{\Re}{\mathop{\mathrm{Re}}\nolimits}

% =================================





% ---------------------------
% new definition macro
% ---------------------------
% 便利なマクロ集です

% 内積のマクロ
%   例えば \inner<\pphi | \psi> のように用いる
\def\inner<#1>{\langle #1 \rangle}

% ================================
% マクロを追加する場合のスペース 

%=================================



% 位相空間
\newcommand{\Int}[1][]{\mathop{\mathrm{Int}}\nolimits_{#1}}
\newcommand{\Cl}[1][]{\mathop{\mathrm{Cl}}\nolimits_{#1}}
\newcommand{\Bd}[1][]{\mathop{\mathrm{Bd}}\nolimits_{#1}}
\newcommand{\bd}{\mathord{\partial}}
\newcommand{\Fr}[1][]{\mathop{\mathrm{Fr}}\nolimits_{#1}}

\newcommand{\dom}[1][]{\mathop{\mathrm{dom}}\nolimits_{#1}}


\newcommand{\mtan}[1][]{T\hspace{-1.75pt}{#1}}
\newcommand{\mtp}[1][]{{}^{t\!}{#1}} % transpose of the morphism
%\newcommand{\mpb}[1][]{{}^{t\!}{#1}'} % pull-back of the morphism
\newcommand{\mpb}[1][]{{#1}_{\pi}} % pull-back of the morphism
\newcommand{\mdiff}[1][]{{#1}_{d}} % pull-back of the morphism
\newcommand{\mptan}[1][]{{#1}_{\tau}} % pull-back of the morphism


%\newcommand{\pb}[1][]{\mathbin{\mathop{\times}\limits_{#1}}}

\newcommand{\cotb}[1][]{T^{\ast}{#1}}
\newcommand{\conm}[2][]{T_{#1}^{\hspace{0.7pt}\ast}{#2}}
\newcommand{\norb}[2][]{T_{#1}{#2}}

\newcommand{\cotz}[1][]{T^{\ast}_{#1}{#1}}

\newcommand{\cotn}[1][]{\dot{T}^{\ast}{#1}}


\newcommand{\muhom}{\mu hom}
\newcommand{\muHom}[1][]{\mathrm{Hom}^\mu_{\raise1.5ex\hbox to.1em{}#1}}

\newcommand{\mucat}[2][]{\mathsf{D}^{\mathrm{b}}_{#1}(#2)}
%\newcommand{\mucat}[2][]{\mathsf{D}^{\mathrm{b}}_{#1}(#2)}

\newcommand{\conv}[1][]{\mathbin{\mathop{\circ}\limits_{#1}}}

\newcommand{\torus}{\mathbf{T}}

\newcommand{\diag}[1][]{\Delta_{#1}}
%  \[\mu_{\Delta_{Y}}\mathrm{R}\mathcal{H}om(G,F)\]
\newcommand{\micro}[1]{\mu_{#1}\hspace{-2pt}}

% ----------------------------
% documenet 
% ----------------------------
% 以下, 本文の執筆スペースです. 
% Your main code must be written between 
% begin document and end document.
% ---------------------------

\begin{document}
\maketitle
\tableofcontents
\section{圏\(\Dompb(X;\varOmega)\)}

\(V\)を\(\cotb[X]\)の部分集合とする．
\(\Dompb(X)\)の充満部分圏\(\mucat[V]{X}\)を次で定める．
\begin{equation}\label{eq:6.1.1}
    \Ob(\mucat[V]{X})\coloneqq
    \left\{F\in\Ob(\Dompb(X));\muS(F)\subset V\right\}
\end{equation}
\(\mucat[V]{X}\)は，
\(\mres[{[1]}]{\mucat[V]{X}}\)をシフト関手，
\[
    \left\{F\to G\to H\to F[1]\mathrel{;}\text{\(\Dompb(X)\) の完全三角で} F,G,H\in\mucat[V]{X}\right\}
\]
を完全三角の族として，三角圏になる．
また，\(\Ob(\mucat[V]{X})\)は\(\Dompb(X)\)の
零系 (null system) をなす．
したがって，\(\Dompb(X)\)を\(\Ob(\mucat[V]{X})\)で局所化できる．

\begin{proof}[{\(\mucat[V]{X}\)が三角圏になることの証明．}]
    まず\(\mucat[V]{X}\)が加法圏であることを示す．

    超局所台はシフトによらないので，
    シフト関手は自己同値\(\mucat[V]{X}\to\mucat[V]{X}\)をひきおこす．
\end{proof}


\clearpage
\begin{leftbar}
\begin{PRP}[{\cite[Prop.6.1.4]{KS90}}]\label{prp.6.1.4}
    \(X\)を多様体とし，\(x_0\in X\)とする．
    \(K\)を\(\conm[x_0]{X}\)の閉固有凸錐とし，
    \(U\)を\(K\)の開部分錐とする．
    \(F\in\Dompb(X)\)とする．
    \(W\)を\(K\cap\muS(F)-\{0\}\)の錐状近傍とする．
    このとき，\(F'\in\Dompb(X)\)と\(\Dompb(X)\)の射\(
        u\colon F'\to F
    \)で次をみたすものが存在する．
    \begin{align}
        \tag{6.1.7}\label{eq:6.1.7}
        \text{\(u\)は\(U\)上で同形である．}\\
        \tag{6.1.8}\label{eq:6.1.8}
        \pi^{-1}(x_0)\cap\muS(F')\subset W\cup\{0\}.
    \end{align}
\end{PRP}
\end{leftbar}

\begin{proof}
    \(W\)が\(K\cap\muS(F)-\{0\}\)の開近傍を含むので，
    必要なら\(K\)を太くして，
    \(\overline{U}\subset\Int{K}\cup\{0\}\)としてよい．
    \(X\)はベクトル空間で\(x_0=0\)としてよい．

    \(\conm[x_0]{X}\)の閉固有凸錐\(\gamma\)を，
    \(K^{\circ a}\subset\gamma U^{\circ a}\)をみたし，
    \(\gamma-\{0\}\)が\(C^1\)級境界\(\delta\gamma\)をもつようにとる．
    \begin{CMT}
        \cite[Thm25.5]{Roc}\begin{quotation}
            \(f\)を\(\rr^n\)上の proper convex function とし，
            \(D\)を\(f\)が微分可能である点の集合とする．
            このとき\(D\)は\(\Int(\dom{f})\)の稠密な部分集合であり，
            \(D\)の\(\Int(\dom{f})\)における補集合は測度\(0\)の集合である．
            さらに，gradient mapping \(
                \nabla{f}\colon x\to\nabla{f}(x)
            \)は\(D\)上で連続である，
        \end{quotation}
        と Sardの定理からいけるか？
    \end{CMT}
    \(X\)の座標\((x)=(x_1,\dots,x_n)\)を
    \begin{equation}
        \tag{6.1.9}\label{eq:6.1.9}
        \gamma\subset\{0\}\cup\left\{x;x_1<0\right\}
    \end{equation}
    をみたすようにとる．
    \(\cotb[X]\)の\((x)\)に対応する座標を\((x;\xi)\)とする．
    \(\gamma^{\circ a}\subset K\)より，\(0\)の開近傍\(\varOmega\)で
    \begin{equation}
        \tag{6.1.10}\label{eq:6.1.10}
        W\supset G=\left\{\xi\in \gamma^{\circ a};\exists(x;\xi)\in\pi^{-1}(\varOmega)\cap \muS(F)\right\}
    \end{equation}
    をみたすものが存在する．
    \begin{CMT}
        ノート参照．
    \end{CMT}
    この\(\varOmega\)に対し，実数\(\varepsilon>0\)を
    \begin{equation}
        \tag{6.1.11}\label{eq:6.1.11}
        \varOmega\supset\left\{x\in\gamma;x_1\geqq-\varepsilon\right\}
    \end{equation}
    となるものをとる．
    このとき，\(X\)の閉集合\(Z\)で次をみたすものが構成できる．
    \begin{align}
        \notag \text{\(Z\)の境界}\delta Z&\text{は\(C^1\)級である．}\\
        \tag{6.1.12}\label{eq:6.1.12}
        0&\in\Int{Z},\\
        \tag{6.1.13}\label{eq:6.1.13}
        Z&\subset\left\{x;x_1\geqq-\varepsilon\right\},\\
        \tag{6.1.14}\label{eq:6.1.14}
        \text{\(\delta Z\)と\(\delta \gamma\)は共通部部で接する．}&\text{すなわち\(\delta Z\cap\delta \gamma\)で}\\
        \notag
        N_x^{\ast}(Z)^{a}&=N_x^{\ast}(\gamma)
    \end{align}
    \begin{CMT}
        構成わからない．
    \end{CMT}
    以上の\(\gamma\)と\(Z\)に対し，\(F'\)を
    \[
        F'=\phi_{\gamma}^{-1}\Rder{{\phi_{\gamma}}_{\ast}}\RG_{Z}(F)
    \]とおき，自然な射\(F'\to F\)を\(u\)とおく．
    この\(F'\)と\(u\)が条件\eqref{eq:6.1.7}, \eqref{eq:6.1.8}をみたすことを示す．
    命題5.2.3\begin{quote}
        \(\phi_{\gamma}^{-1}\Rder{{\phi_{\gamma}}_{\ast}}\RG_{Z}(F)
        \to
        \RG_{Z}(F)\)は\(X\times\Int{\gamma^{\circ a}}\)で同形
    \end{quote}
    であり，\(\RG_{Z}(F)\to F\)は\(\Int{Z}\)で同形なので
    \(u\)は\(\Int{Z}\times\Int{\gamma^{\circ a}}\)で同形である．
    いま，\(U\subset\Int{\gamma^{\circ a}}\)なので\eqref{eq:6.1.7}がなりたつ．

    再び命題5.2.3 (i) (\(\Leftarrow\)) から
    \[
        \muS(F')\subset X\times\gamma^{\circ a}
    \]
    なので，あと示すべきは
    \begin{equation}
        \tag{6.1.15}\label{eq:6.1.15}
        \begin{cases*}
            \xi\in\delta(\gamma^{\circ a})-\{0\},\\
            (0;\xi)\in\muS(F')
        \end{cases*}
        \Longrightarrow
        \xi\in G
    \end{equation}
    となること．（ノート）

    \(q_i\colon X\times X\to X\)を射影，\(s\colon X\times X\to X\)を\(s(x,x')=x-x'\)とする．
    命題3.5.4\begin{quotation}
        \(F\in\Dompb(X)\)とする．\(Z(\gamma)=\left\{(x,y)\in X\times X;y-x\in\gamma\right\}\)とする．
        このとき
        \[
            \Rder{{q_1}_{\ast}}\left(\left(q_2^{-1}F\right)_{Z(\gamma)}\right)\cong
            \phi_{\gamma}^{-1}\Rder{{\phi_{\gamma}}_{\ast}}F
        \]である．
    \end{quotation}
    より，
    \begin{equation}\tag{6.1.16}\label{eq:6.1.16}
        F'\cong
        \Rder{{q_2}_{\ast}}\left(s^{-1}A_{\gamma}\tens[]q_1^{-1}\RG_Z(F)\right)
    \end{equation}
    がなりたつ．
    条件\eqref{eq:6.1.9}, \eqref{eq:6.1.13}より，
    \(q_2:s^{-1}(\gamma)\cap q_1^{-1}(Z)\to X\)は固有である．（ノート）

    よって，SSの評価式が使える．
    \begin{align*}
        \muS(F')&=\muS(
        \Rder{{q_2}_{\ast}}\left(s^{-1}A_{\gamma}\tens[]q_1^{-1}\RG_Z(F)\right)
        )\\
        &\subset
        \mpb[q_2]\mdiff[q_2]^{-1}\muS(s^{-1}A_{\gamma}\tens[]q_1^{-1}\RG_Z(F))\\
        &\subset
        \mpb[q_2]\mdiff[q_2]^{-1}\muS(s^{-1}A_{\gamma})+\muS(q_1^{-1}\RG_Z(F))\\
        &\subset
        \mpb[q_2]\mdiff[q_2]^{-1}\mdiff[s]\mpb[s]^{-1}\muS(A_{\gamma})+\mdiff[q_1]\mpb[q_1]^{-1}\muS(\RG_Z(F)).
    \end{align*}
    \((0;\xi)\in\muS(F')\)に対し，
\end{proof}




















%\section{余接束内の法錐}
%\section{順像}
%\section{超局所化}
%\section{包合性と伝播}
%\section{包合的多様体の近傍における層}
%\section{超局所化と逆像}

\clearpage
\appendix
%\section{ノート中で用いる主張}

\section{三角圏}\label{ssec:tricat}

\begin{DFN}
    加法圏とは次の3つの
    条件(\ref{additive:bilin})--(\ref{additive:biproduct})を
    みたす圏のことである．
    \begin{enumerate}[(1)]
        %\renewcommand{\labelenumi}{({\arabic{enumi}})}
        \item \(\cat\)の任意の対象の組\((X, Y)\)に対し，
        \(\Hom_{\cat}(X,Y)\)が加法群になり，
        %どの対象$X, Y, Z\in\cat$に対しても
        合成
        %\(
        %    \circ\colon
        %    \Hom_{\cat}(Y,Z)\times\Hom_{\cat}(X,Y)
        %    \to
        %    \Hom_{\cat}(X,Z)
        %\)
        が双線型である．\label{additive:bilin}
        \item 零対象$0\in\cat$が存在する．
        さらに$\Hom_{\cat}(0,0)=0$が成り立つ．\label{additive:zero}
        \item 任意の対象$X, Y\in\cat$に対して積と余積が存在し，
        さらにそれらは同型になる．（それらを複積といい$X\oplus Y$とかく．）\label{additive:biproduct}
    \end{enumerate}
\end{DFN}

$\cat$を加法圏とし，$T\colon\cat\to\cat$を自己関手とする．
$\cat$の三角とは射の列
\begin{equation*}
    X\to Y\to Z\to TX
\end{equation*}
のことである．

\begin{Definition}
    三角圏$\cat$は次のデータ\eqref{eq:tricat}, \eqref{eq:family-dt}と
    規則(TR\ref{TR0})--(TR\ref{TR5})からなる．
    \begin{description}%\label{des:data}
        \item[データ] \begin{align}
            %\renewcommand{\labelenumi}{(\ref{ssec:tricat}.{\arabic{enumi}})}
            \text{加法圏}\cat\text{と自己同値}T\colon\cat\to\cat\text{の組},\label{eq:tricat}\\
            \text{特三角 (distinguished triangle) の族．}\label{eq:family-dt}
        \end{align}    
        \item[規則] \label{des:data}
        \begin{quote}
            \begin{enumerate}
                %\setlength{\leftskip}{11pt}
                \setcounter{enumi}{-1}
                \renewcommand{\labelenumi}{(TR{\arabic{enumi}})}
                \item 特三角に同形な三角は特三角である．\label{TR0}
                \item 任意の対象$X\in\cat$に対し，
                    $X\overset{\id_X}{\longrightarrow}X\longrightarrow
                    0\longrightarrow TX$は特三角である．\label{TR1}
                \item $\cat$の任意の射$f\colon X\to Y$は
                    特三角$X\overset{f}{\to}Y\to Z\to TX$に埋め込める．
                    つまり$Z\in\cat$で$X\overset{f}{\to}Y\to Z\to TX$が
                    特三角となるものが存在する．\label{TR2}
                \item $X\overset{f}{\to}
                    Y\overset{g}{\to} 
                    Z\overset{h}{\to} TX$が
                    特三角であることと$Y\overset{g}{\longrightarrow}
                    Z\overset{h}{\longrightarrow} TX\overset{-T(f)}{\longrightarrow} TY$が
                    特三角であることは同値である．\label{TR3}
                \item ２つの特三角$X\overset{f}{\to}
                Y\to Z\to TX$, $X'\overset{f'}{\to}
                Y'\to Z'\to TX'$に対し，可換図式
                \begin{equation*}
                    \vcenter{\xymatrix
                    @C=26pt@R=26pt
                    {
                    X\ar[r]^-{f}\ar[d]^-{u}
                    &
                    Y\ar[d]^-{v} 
                    \\
                    X'\ar[r]^-{f'}
                    &
                    Y'
                    }}
                \end{equation*}
                は特三角の射に埋め込める．\label{TR4}
                \item\label{TR5}（八面体公理）．３つの特三角
                \begin{align*}
                    X\overset{f}{\to}Y\overset{h}{\to} Z'\to TX, \\
                    Y\overset{g}{\to}Z\overset{k}{\to} X'\to TY,\\
                    X\overset{g\circ f}{\longrightarrow}
                    Z\overset{l}{\to} Y'\to TX
                \end{align*}
                に対し，完全三角\(
                    Z'\overset{u}{\to}
                    Y'\overset{v}{\to} 
                    X'\overset{w}{\to} 
                    TZ'
                \)で以下の図式を可換にするものが存在する．
                \[\vcenter{\xymatrix{
                    X
                    \ar[r]^-{f} 
                    \ar[d]^-{\id}
                    & 
                    Y\ar[r]^-{h} 
                    \ar[d]^-{g}
                    & 
                    Z'\ar[r]^-{}
                    \ar@{..>}[d]^-{u}
                    & 
                    TX\ar[d]^-{\id}\\
                    X\ar[r]^-{g\circ f} 
                    \ar[d]^-{f}
                    & 
                    Z\ar[r]^-{l} 
                    \ar[d]^-{\id}
                    & 
                    Y'\ar[r]^-{}
                    \ar@{..>}[d]^-{v}
                    & 
                    TX\ar[d]^-{T(f)}\\
                    Y\ar[r]^-{g} 
                    \ar[d]^-{h}
                    & 
                    Z\ar[r]^-{k} 
                    \ar[d]^-{l}
                    & 
                    X'\ar[r]^-{}
                    \ar[d]^-{\id}
                    & 
                    TY\ar[d]^-{T(h)}\\
                    Z'
                    \ar@{..>}[r]^-{u} 
                    & 
                    Y'\ar@{..>}[r]^-{v} 
                    & 
                    X'\ar@{..>}[r]^-{w}
                    & 
                    TZ'.
                }}\]
            \end{enumerate}    
        \end{quote}
    \end{description}
\end{Definition}


\begin{leftbar}
\begin{PRP}[{\cite[Prop.3.4.4.]{KS90}}]\label{prp3.4.4}
  \(X\)と\(Y\)を有限c柔軟次元局所コンパクト（ハウスドルフ）空間とし，
  \(q_1\)と\(q_2\)をそれぞれ\(X\times Y\)から\(X\)と\(Y\)への射影とする．
  \(F\in\Dompb(A_X)\), \(G\in\Dompb(A_Y)\)とし，\(F\)はコホモロジー構成可能層とする．
  このとき，
  \[
    \vD{F}\letens G\to \RHOM(q_1^{-1}F,q_2^{!}G)
  \]は同型である．
  \(X\)が\(C^0\)多様体ならば，\[
    \vD'{F}\letens G\to \RHOM(q_1^{-1}F,q_2^{-1}G)
  \]も同型である．
\end{PRP}
\end{leftbar}

\begin{leftbar}
\begin{PRP}[{\cite[Prop.4.3.4]{KS90}}]\label{prp4.3.4}
  \(G\in\Dompb(Y)\)とする．
  このとき，標準的な射による可換図式    
  \[    
      \vcenter{\xymatrix@C=36pt@R=36pt{
      \Rder{{\mpb[f_{N}]}_{!}}{\mdiff[f_{N}]^{-1}}\mu_{N}(G)
      \ar[d]
      \ar[r]
      %\ar@{}[dr]|\square
      &
      \mu_{M}(\Rder{f_{!}}G)
      %\ar[r]
      \ar[d]
      \\
      \Rder{{\mpb[f_{N}]}_{\ast}}({\mdiff[f_{N}]^{!}}\mu_{N}(G)\tens[]\omega_{Y/X}\tens[]\omega_{N/M}^{\otimes{-1}})
      %\ar@{}[dr]|\square
      &
      \mu_{M}(\Rder{f_{\ast}}G)
      \ar[l]
      }}
  \]が存在する．
  さらに，\(\supp(G)\to X\)と\(
      C_N(\supp(G))\to T_{M}X
  \)が固有かつ\(\supp(G)\cap f^{-1}(M)\subset N\)ならば，
  以上の射はすべて同型である．
  
  とくに\(f^{-1}(M)=N\)かつ，
  \(f\)が\(M\)に関して鮮明\footnotemark[2] (clean) かつ\(\supp(G)\)上固有のとき，
  以上の射はすべて同型である．
\end{PRP}
\end{leftbar}

\footnotetext[2]{試験的な訳語}

\begin{leftbar}
\begin{PRP}[{\cite[Prop.4.3.5]{KS90}}]\label{prp4.3.5}
  \(F\in\Dompb(X)\)とする．
  \begin{enumerate}[(i)]
      \item 標準的な射で次の図式を可換にするものが存在する．\label{prp4.3.5.1}
      \[\vcenter{\xymatrix@C=56pt@R=46pt{
          \Rder{{\mdiff[f_N]}_{!}}(\omega_{N/M}\tens[]{\mpb[f_N]^{-1}}\mu_{M}(F))
          \ar[d]^{}
          \ar[r]_{}
          %\ar@{}[dr]|\square
          &
          \mu_{N}(\omega_{Y/X}\tens[]{f^{-1}}F)
          %\ar[r]
          \ar[d]
          \\
          \Rder{{\mdiff[f_N]}_{\ast}}{\mpb[f_N]^{!}}\mu_{M}(F)
          %\ar@{}[dr]|\square
          &
          \mu_{N}({f^{!}}F)
          \ar[l]_{\beta}
      }}\]
      ここで縦向きの矢印は(3.1.6)から導かれるものである．
      \item \label{prp4.3.5.2}
      \(f\colon Y\to X\)と\(\mres[f]{N}\colon N\to M\)が
      平滑ならば，これらの射はすべて同型となる．    
      \item \label{prp4.3.5.3}
      \(f\)が\(M\)に横断的かつ\(f^{-1}(M)=N\)ならば，
      自然な射\[
          \Rder{{\mdiff[f_N]}_{\ast}}\mpb[f_N]^{-1}\mu_{M}(F)
          \to \mu_{N}(f^{-1}F)
      \]が存在する．
  \end{enumerate}
\end{PRP}
\end{leftbar}














%\section{aa}
%===============================================
% 参考文献スペース
%===============================================
\begin{thebibliography}{20} 
  \bibitem[dR]{dR} 
  de Rham, \textit{Vari\'et\'es diff\'erentiables}, Hermann, 1953.
  和訳：ド・ラーム, 微分多様体, 東京図書, 1974.
  \bibitem[GKS12]{GKS12} St\'ephane Guillermou, 
  Masaki Kashiwara, Pierre Schapira, 
  \textit{Sheaf Quantization of Hamiltonian Isotopies and Applications to Nondisplaceability Problems}, 
  Duke. Math. J. 161, No. 2, pp.201--245, 2012.  
  \bibitem[GP74]{GP74} Victor Guillemin, Alan Pollack, 
    \textit{Differential Topology}, 
    Prentice-Hall, 1974.
  \bibitem[KS90]{KS90} Masaki Kashiwara, Pierre Schapira, 
    \textit{Sheaves on Manifolds}, 
    Grundlehren der Mathematischen Wissenschaften, 292, Springer, 1990.
  \bibitem[Hat89]{Hat89} 服部晶夫, 多様体 増補版, 岩波全書 288, 岩波書店, 1989.
  %\bibitem[Og02]{Og02} 小木曽啓示, 代数曲線論, 朝倉書店, 2022.
  \bibitem[Hor63]{Hor63} Lars H\"ormander, 
  \textit{Linear Pertial Differential Operators}, Fourth Printing, 
  Grundlehren der Mathematischen Wissenschaften, 116, Springer, 1963.
  \bibitem[Hor71]{Hor71} Lars H\"ormander, 
  \textit{Fourier Integral Operators I}, 
  Acta Math. 127, 79--183, (1971).
  \bibitem[Hor89]{Hor89} Lars H\"ormander, 
  \textit{The Analysis of Linear Pertial Differential Operators I}, Second Edition,
  Grundlehren der Mathematischen Wissenschaften, 256, Springer, 1989.
  \bibitem[Hor85]{Hor85} Lars H\"ormander, 
  \textit{The Analysis of Linear Pertial Differential Operators III}, 
  Grundlehren der Mathematischen Wissenschaften, 274, Springer, 1985.
  \bibitem[Roc]{Roc} Rockaffeller, \textit{Convex Analysis},. 
  \bibitem[Sch85]{Sch85} Pierre Schapira, 
    \textit{Microdifferential Systems in the Complex Domain}, 
    Grundlehren der Mathematischen Wissenschaften, 269, 
    Springer, 1985.
    \bibitem[ST92]{ST92} Pierre Schapira, Nobuyuki Tose, 
    \textit{Morse Inequalities for \(\rr\)-constructible Sheaves}, 
    Adv. Math. 93, pp.1--8, 1992. 
    \bibitem[Zwo12]{Zwo12} Maciej Zworski,
    \textit{Semiclsaaical Analysis}, 
    Graduate Studies in Mathematics, 138, 
    American Mathematical Society, 
    2012. 
\end{thebibliography}

%===============================================




\end{document}
